\section{Discussion}
\label{sec:discussion}

\subsection{Key Findings and Their Implications}

\textbf{Finding 1: Both HF Hub and OpenML exhibit domain-temporal selection bias, not random coverage gaps.}

The 30\% HF coverage is concentrated in NLP benchmarks and small-scale image classification benchmarks (CIFAR-10, CIFAR-100, SVHN) that have been adopted by the HF ecosystem. The 22\% OpenML coverage is predominantly classical tabular datasets (iris, wine, adult, covertype, diabetes), with NLP benchmarks (IMDB, 20newsgroups) and MNIST also present. Large-scale DL vision (ImageNet, COCO, Pascal VOC, KITTI), speech (LibriSpeech, TIMIT, WSJ), and graph (Cora, Citeseer) datasets --- which dominate Raff's post-2012 deep learning papers --- are absent from both platforms.

\textbf{Finding 2: Papers With Code is the appropriate primary data source for a redesigned study.}

Papers With Code explicitly tracks paper-dataset linkages across all dataset domains and historical periods. Unlike HF Hub (NLP-centric) and OpenML (tabular-centric), PwC is designed for the paper-dataset-code linkage we require and covers Raff's corpus temporal period. The theoretical hypothesis --- documentation completeness and concentration predict reproducibility failure --- remains scientifically plausible and worth testing, but requires PwC as the data source.

\textbf{Finding 3: Retroactively-created HF cards are systematically thin.}

The 0.41 mean field-presence score reveals that even when a card exists, it may not contain the fields most relevant to the misuse mechanism. The \texttt{intended\_use} and \texttt{out\_of\_scope\_use} fields --- which directly encode scope boundaries --- are systematically absent from retroactively-created cards, consistent with the timing of the Datasheets schema publication.

\subsection{Limitations}

\textbf{Limitation 1: The original research question remains unanswered.} The central hypothesis --- do metadata signals predict reproducibility failure? --- is neither confirmed nor refuted. We report an infrastructure characterization. Running regression on 30\% IV1 coverage would produce biased, uninterpretable results; blocking was the correct decision.

\textbf{Limitation 2: H-E1 used canonical dataset names only.} Alternate name variants (e.g., ``uoft-cs/cifar10'') might recover additional HF cards. However, a scalable auditing pipeline requiring manual name variant curation defeats the purpose of automated auditing; the DL domain gap is structural in any case.

\textbf{Limitation 3: Results are specific to the Raff 2019 corpus.} Post-2019 papers with more NLP and HF-native datasets would show substantially higher HF coverage.

\textbf{Limitation 4: The theoretical mechanism remains empirically unverified.} The three-step causal chain (concentration $\to$ artifact accumulation; incompleteness $\to$ out-of-context application; joint effect $\to$ reproducibility failure) is theoretically grounded in \citet{DAm2021Underspecification} and \citet{Gebru2021Datasheets} but not yet empirically confirmed.

\subsection{Broader Impact}

Our infrastructure characterization directly benefits researchers planning HF+OpenML-based reproducibility auditing of pre-2018 ML literature. Without this characterization, such researchers would build pipelines that silently fail to represent 70--78\% of the target corpus. The validated pipeline components are immediately reusable for the redesigned PwC-based study.

A reflexivity note: the Datasheets schema we use to measure completeness was designed for new dataset creation, not retroactive annotation of 1990s benchmarks. Evaluating CIFAR-10 by whether its HF card contains an \texttt{out\_of\_scope\_use} field conflates absence-of-documentation with absence-of-concern-for-misuse. A redesigned study should address this temporal validity problem in the IV operationalization.
